% !TeX document-id = {f19fb972-db1f-447e-9d78-531139c30778}
% !BIB program = biber
\documentclass[compress]{beamer}
\usepackage[T1]{fontenc}
\usetheme[block=fill,subsectionpage=progressbar,sectionpage=progressbar]{metropolis} 

\usepackage{wasysym}
\usepackage{etoolbox}
\usepackage[utf8]{inputenc}

\usepackage{threeparttable}
\usepackage{subcaption}

\usepackage{tikz-qtree}
\setbeamercovered{still covered={\opaqueness<1->{5}},again covered={\opaqueness<1->{100}}}


\usepackage{listings}

\lstset{
	basicstyle=\scriptsize\ttfamily,
	columns=flexible,
	breaklines=true,
	numbers=left,
	%stepsize=1,
	numberstyle=\tiny,
	backgroundcolor=\color[rgb]{0.85,0.90,1}
}

 
\usepackage[american]{babel}
\usepackage{csquotes}
\usepackage[style=apa, backend = biber]{biblatex}
\DeclareLanguageMapping{american}{american-UoN}
\addbibresource{../../bdaca/bdaca.bib}
\graphicspath{{../../bdaca/pictures/}}
\renewcommand*{\bibfont}{\tiny}

\usepackage{tikz}
\usetikzlibrary{shapes,arrows,matrix}
\usepackage{multicol}

\usepackage{subcaption}

\usepackage{booktabs}
\usepackage{graphicx}


 
\makeatletter
\setbeamertemplate{headline}{%
  \begin{beamercolorbox}[colsep=1.5pt]{upper separation line head}
  \end{beamercolorbox}
  \begin{beamercolorbox}{section in head/foot}
    \vskip2pt\insertnavigation{\paperwidth}\vskip2pt
  \end{beamercolorbox}%
  \begin{beamercolorbox}[colsep=1.5pt]{lower separation line head}
  \end{beamercolorbox}
}
\makeatother

 

\setbeamercolor{section in head/foot}{fg=normal text.bg, bg=structure.fg}



\newcommand{\question}[1]{
	\begin{frame}[plain]
		\begin{columns}
			\column{.3\textwidth}
			\makebox[\columnwidth]{
				\includegraphics[width=\columnwidth,height=\paperheight,keepaspectratio]{mannetje.png}}
			\column{.7\textwidth}
			\large
			\textcolor{orange}{\textbf{\emph{#1}}}
		\end{columns}
\end{frame}}




\title[Big Data and Automated Content Analysis]{\textbf{Teaching the Teacher: Python} \\ Day 3 -- Morning: »Data Formats«}
\author[Damian Trilling]{Damian Trilling \\ ~ \\ \footnotesize{d.c.trilling@uva.nl \\@damian0604} \\ \url{www.damiantrilling.net}}
\date{30 June 2021}
\institute[UvA]{Afdeling Communicatiewetenschap \\Universiteit van Amsterdam}

\begin{document}

\begin{frame}{}
\titlepage
\end{frame}

\begin{frame}{Today}
\tableofcontents
\end{frame}


\begin{frame}[standout]
Everything clear from yesterday?
\end{frame}





\section{Data structures and files}

\begin{frame}[standout]
	(a bit more in-depth than yesterday)
\end{frame}

\begin{frame}{Some ways of storing data}
\begin{table}[]
\resizebox{\textwidth}{!}{%
\begin{tabular}{@{}llll@{}}
\toprule
use case          & data type       & structure         & typical file format \\ \midrule
\onslide<+->{(long) texts       & string          & unstructured      & multiple .txt files               \\}
\onslide<+->{list of messages/words/\ldots       & list of strings  & newline-delimited      & .txt                \\}
\onslide<+->{hierarchical data & dict                          & (semi-)structured & .json               \\}
\onslide<+->{table             & nested lists, dict, dataframe & tabular           & .csv (.json)   \\ \bottomrule}
\end{tabular}%
}
\end{table}
\pause
\textcolor{red}{\footnotesize{(Of course, there are many, many, many other formats we can use)}}
\end{frame}




\begin{frame}[fragile]{string $\leftrightarrow$ file}
\begin{lstlisting}
data = "Hi I'm a string."
with open("mytext.txt", mode="w") as fo:
    fo.write(data)
\end{lstlisting}
$\Rightarrow$ create (or overwrite(!)) file, assign handler name \texttt{fo}, write string to it.

\pause 
\begin{lstlisting}
with open("mytext.txt", mode="r") as fi:
    data = fi.read()
\end{lstlisting}

$\Rightarrow$ make connection with file for reading, assign handler name \texttt{fi}, read string from it

\pause

%\textcolor{red}{Of course, you can also loop over \emph{multiple]} files: e.g, read all files from a dictionary and put them into a list $\Rightarrow$ just nest these commands in a for loop!}
\end{frame}

\question{For what can you use this?}


\begin{frame}[fragile]{list $\rightarrow$ file}
\begin{lstlisting}
data = ["ik", "jij", "je", "hij", "zij", "ze", "wij", "we", "jullie"]

with open("pronouns.txt", mode="w") as fo:
    for pronoun in data:
        fo.write(pronoun)
        fo.write("\n")
\end{lstlisting}
$\Rightarrow$ create file, assign handler name \texttt{fo}, write each element from list to \texttt{fo} followed by a line break

Result: a file \texttt{pronouns.txt} with this content:
\begin{lstlisting}
ik
jij
je
hij
zij
ze
wij
we
jullie
\end{lstlisting}
\end{frame}



\begin{frame}[fragile]{list $\leftarrow$ file}
\begin{lstlisting}
with open("pronouns.txt", mode="r") as fi:
     data = [line.strip() for line in fi]
print(data)
\end{lstlisting}
$\Rightarrow$ Open file for reading, assign handler name \texttt{fi}, loop over all lines in \texttt{fi}, strip whitespace from end (such as line endings), put in new list.

Output:
\begin{lstlisting}
['ik', 'jij', 'je', 'hij', 'zij', 'ze', 'wij', 'we', 'jullie']
\end{lstlisting}
\end{frame}


\question{For what can you use this?}





\begin{frame}[fragile]{dict $\rightarrow$ file}
\begin{lstlisting}
import json

data = {'Alice': {'office': '020222', 'mobile': '0666666'},
    'Bob': {'office': '030111'},
    'Carol': {'office': '040444', 'mobile': '0644444'},
    'Daan': "020222222",
    'Els': ["010111", "06222"]}

with open("phonebook.json", mode="w") as f:
    json.dump(data, f)

\end{lstlisting}
$\Rightarrow$ Open file for writing, convert dict to JSON, dump in file.

Creates a file \texttt{phonebook.json} that looks like this:
\begin{lstlisting}
{"Alice": {"office": "020222", "mobile": "0666666"}, "Bob": {"office": "030111"}, "Carol": {"office": "040444", "mobile": "0644444"}, "Daan": "020222222", "Els": ["010111", "06222"]}
\end{lstlisting}
\end{frame}




\begin{frame}[fragile]{dict $\leftarrow$ file}
\begin{lstlisting}
import json

with open("phonebook.json", mode="r") as f:
	data = json.load(f)
	
\end{lstlisting}

$\Rightarrow$ Reads data from f, converts to dict (or list of dicts\ldots), store in \texttt{data}

\pause

\small
\begin{alertblock}{JSON lines}
There is also a dialect in which you write \emph{one JSON object per line instead of per file}. For this, you can use a for loop as in the one-string-per-file example, but convert each string with  \texttt{json.load\textbf{s}} or \texttt{json.dump\textbf{s}} to a dict.
\end{alertblock}


\end{frame}


\question{For what can you use this?}




\begin{frame}[fragile]{Tabular data}
How can we store this data?
\begin{lstlisting}
names = ['Alice', 'Bob', 'Carol', 'Daan', 'Els']
phonenumbers = ['020111111', '020222222', '020333333', '020444444', '020555555']
\end{lstlisting}
\pause

1. We could convert to some dict and store as json (not too bad\ldots)

\pause
2. We can store in a table
(next slide)
\end{frame}

\begin{frame}[fragile]{Tabular data}
\begin{lstlisting}
import csv
with open("mytable.csv", mode="w") as f:
    mywriter = csv.writer(f)
    for row in zip(names,phonenumbers):
        mywriter.writerow(row)
\end{lstlisting}
	
Results in a file \texttt{mytable.csv} that looks like this:
\begin{lstlisting}
Alice,020111111
Bob,020222222
Carol,020333333
Daan,020444444
Els,020555555
\end{lstlisting}
\pause
\textcolor{red}{But you don't have to do it like this! Often, pandas is more user-friendly. Yet, also pandas doesn't save you from thinking about encodings and dialects\ldots}
\end{frame}


\section{Encodings and dialects}


\begin{frame}{Choices to make}
For all text-based (as opposed to binary) file formats:
\begin{block}{How to separate data?}
	\begin{itemize}
		\item new line = new record?
		\item Unix-style (\texttt{\textbackslash n}, also known as LF),  or Windows-style (\texttt{\textbackslash r\textbackslash n}, also known as CRLF) line endings?
		\item some delimiter = new field?
		\item or new file = new record?
	\end{itemize}
\end{block}
\pause
\begin{block}{How to convert bytes to characters?}
	\begin{itemize}
		\item choose right encoding (e.g., UTF-8)
		\item (seldom) does the file start with a so-called byte-order-marker (BOM) -- then the encoding is often referred to as sth like UTF-8-BOM
	\end{itemize}
\end{block}
\end{frame}


\begin{frame}{Let's look at csv files}
	\begin{block}{comma-separated values: always a good choice}
		\begin{itemize}
			\item \emph{All} programs can read it
			\item Even human-readable in a simple text editor
			\item Plain text, with a comma (or a semicolon) denoting column breaks
			\item No limits regarging the size
			\item But: several dialects (e.g., \texttt{,} vs. \texttt{;} as delimiter)
			\item Also: tab-separated files (.csv, .tab, .txt -- no consensus) (delimiter is \texttt{\textbackslash t})
		\end{itemize}
	\end{block}
\end{frame}


\begin{frame}[fragile]{A CSV-file with tweets}
	\begin{itemize}
		\item delimiter is \texttt{,}
		\item with  header row
	\end{itemize}
	\begin{lstlisting}
text,to_user_id,from_user,id,from_user_id,iso_language_code,source,profile_image_url,geo_type,geo_coordinates_0,geo_coordinates_1,created_at,time
:-) #Lectrr #wereldleiders #uitspraken #Wikileaks #klimaattop http://t.co/Udjpk48EIB,,henklbr,407085917011079169,118374840,nl,web,http://pbs.twimg.com/profile_images/378800000673845195/b47785b1595e6a1c63b93e463f3d0ccc_normal.jpeg,,0,0,Sun Dec 01 09:57:00 +0000 2013,1385891820
Wat zijn de resulaten vd #klimaattop in #Warschau waard? @EP_Environment ontmoet voorzitter klimaattop @MarcinKorolec http://t.co/4Lmiaopf60,,Europarl_NL,406058792573730816,37623918,en,<a href="http://www.hootsuite.com" rel="nofollow">HootSuite</a>,http://pbs.twimg.com/profile_images/2943831271/b6631b23a86502fae808ca3efde23d0d_normal.png,,0,0,Thu Nov 28 13:55:35 +0000 2013,1385646935
\end{lstlisting}
\end{frame}



\begin{frame}[plain]
	\makebox[\linewidth]{
		\includegraphics[width=\columnwidth,height=.8\paperheight,keepaspectratio]{csv-atom-encoding.png}}

Opening a file in a (good) text editor (here: Atom) to check its encoding and line-ending style.
\end{frame}


\begin{frame}{Why encodings are really, really, REALLY important -- and why you shouldn't let Excel mess with them}
\begin{itemize}
	\item Unicode is around for decades, but sometimes legacy encodings (ASCII, ANSI, Windows-1252, \ldots) are still used
	\item Problem 1: They don't support all Unicode symbols (emoticons, different scripts, \ldots)
	\item Problem 2: What is an \texttt{\"a} in the one encoding may be an \texttt{\o}  in another $\Rightarrow$ big confusion if you use the wrong one
	\item Some programs (looking at you, Excel!) may use legacy encodings when saving CSV files without telling you!
\end{itemize}
\textbf{Make sure to use UTF-8 from beginning to end, unless you know what you are doing!}

\end{frame}



\section{Dataframes}


\begin{frame}{What are dataframes?}
\texttt{pd.DataFrame}s (from the pandas package) are
	\begin{itemize}
		\item objects that store tabular data in rows and columns.
		\item columns and rows can have names
		\item they have methods built-in for data wrangling and analysis
	\end{itemize}

\end{frame}

\begin{frame}{Creating dataframes}
	\begin{columns}[T]
		\column{.5\textwidth}
\textbf{Option 1: transform existing data into a dataframe}
	
\texttt{df = pd.DataFrame(list-of-lists, dict, or similar)}

(use \texttt{pd.DataFrame?} for help if necessary)

		\column{.5\textwidth}
\textbf{Option 2: read from file}
		\makebox[\columnwidth]{
	\includegraphics[width=\columnwidth,height=\paperheight,keepaspectratio]{pd-read.png}}
Using tab-completion to see methods to read dataframes from a file)
	\end{columns}
\end{frame}

\begin{frame}[plain]
	\includegraphics[width=\columnwidth,height=\paperheight,keepaspectratio]{pd-read-csv.png}
\end{frame}



\question{Can you think of a situation when you would use the for-loop approach to reading or writing tabular data instead of the pandas approach? What are pros and cons?}






\begin{frame}{When (not) to use dataframes}
	\footnotesize
\begin{columns}[T]
\column{.5\textwidth}
\begin{exampleblock}{use it!}
\begin{itemize}
	\item tabular data
	\item visual inspection
	\item data wrangling or statistical analysis
\end{itemize}
\end{exampleblock}
\column{.5\textwidth}
\begin{alertblock}{don't use it}
\begin{itemize}
	\item non-tabular data
	\item when it does not make sense to consider rows as ``cases'' and columns as ``variables''
	\item if you only care about one (or maybe two) column anyway
	\item size of dataset $>$ available RAM 
	\item long or expensive operations, play safe and write to / read from file line by line*
\end{itemize}
\end{alertblock}
\end{columns}

\tiny{* imagine scraping 10,000 pages for a week and your program crashes at nr. 9,999 just before saving the dataframe\ldots}
\end{frame}









\section{Beyond standard data files}
\subsection{API's}

\begin{frame}{Beyond files}
It probably has occured to you by now that


\textcolor{red}{\large{\textbf{it's all about the structure}}}:
\begin{itemize}
	\item<2-> we can write \textit{anything} to files
	\item<3-> it doesn't really matter whether sth is called "csv" or whatever -- as long as we know how records are delimited and what the encoding is
	\item<4-> Maybe check out Chapter 9 in the old book for an example of how we can write files in a strange format called GDF (for network data) \emph{even though it is not natively supported} \parencite{Trilling2016} (\url{https://github.com/damian0604/bdaca/tree/master/book})
	\item<5-> and $\ldots$ do we then even need files?
\end{itemize}

\end{frame}

\begin{frame}{APIs}
	{\huge{No.}}
	Why not just send a JSON object (with a question/request) directly through the internet and get another one (with an answer/response) back?
	
	That's what (most) Application Programming Interfaces (APIs) do.
	
	\begin{block}{Great if we have a nested data structure}
		\begin{itemize}
			\item<2-> Items within news feeds
			\item<3-> Personal data within authors within books
			\item<4-> Bio statements within authors within tweets
		\end{itemize}
	\end{block}
\end{frame}


\begin{frame}[fragile]{A JSON object containing GoogleBooks data}
	\begin{lstlisting}
{'totalItems': 574, 'items': [{'kind': 'books#volume', 'volumeInfo': {'publisher': '"O\'Reilly Media, Inc."', 'description': u'Get a comprehensive, in-depth introduction to the core Python language with this hands-on book. Based on author Mark Lutz\u2019s popular training course, this updated fifth edition will help you quickly write efficient, high-quality code with Python. It\u2019s an ideal way to begin, whether you\u2019re new to programming or a professional developer versed in other languages. Complete with quizzes, exercises, and helpful illustrations, this easy-to-follow, self-paced tutorial gets you started with both Python 2.7 and 3.3\u2014 the
...
...
'kind': 'books#volumes'}
	\end{lstlisting}
\end{frame}



\begin{frame}{Who offers APIs?}
	The usual suspects: Twitter, Facebook, Google -- but also Reddit, Youtube, \ldots \\~
	\begin{block}{If you ever leave your bag on a bus on Chicago}<2->
		\onslide<3->{
			\ldots but do have Python on your laptop, watch this: \url{https://www.youtube.com/watch?v=RrPZza_vZ3w}. \\
			That guy queries the Chicago bus company's API to calculate when \emph{exactly the vehicle} with his bag arrives the next time at the bus stop in front of his office.
			\\~\\ \tiny{(Yes, he tried calling the help desk before, but they didn't know. He got his bag back.)}
		}
	\end{block}
\end{frame}

\begin{frame}{APIs}
	\begin{block}{Pro}<1->
		\begin{itemize}
			\item Structured data (JSON!)
			\item Easy to process automatically
			\item Can be directy embedded in your script
		\end{itemize}
	\end{block}
	\begin{block}{Con}<2->
		\begin{itemize}
			\item Often limitations (requests per minute, sampling, \ldots)
			\item You have to trust the provider that he delivers the right content \parencite{Morstatter2013} 
			\item {\textbf{Some APIs won't allow you to go back in time!}}
		\end{itemize}
	\end{block}


$\Rightarrow$ We work with a simple API in the exercise.

$\Rightarrow$ More about APIs versus webscraping this afternoon
\end{frame}


\subsection{Messy data}

\begin{frame}
	Collecting data:\\
	Parsing text files
\end{frame}

\begin{frame}{For messy input data or for semi-structured data}
	Guiding question: Can we identify some kind of pattern?
	\begin{block}{Examples}<2->
		\begin{itemize}
			\item<3-> LexisNexis gives you a chunk of text (rather than, e.g., a structured JSON object)
			\item<3-> But as long as you can find any pattern or structure in it, you can try to write a Python script to \emph{parse} the data (and put it in a dict, lists, or a dataframe)
		\end{itemize}
	\end{block}
\end{frame}

{\setbeamercolor{background canvas}{bg=black}
	\begin{frame}[plain]
		\makebox[\linewidth]{
			\includegraphics[width=\paperwidth,height=\paperheight,keepaspectratio]{lexisnexistxt.png}
		}
	\end{frame}
	\begin{frame}[plain]
		\makebox[\linewidth]{
			\includegraphics[width=\paperwidth,height=\paperheight,keepaspectratio]{lexisnexistxt2.png}
		}
	\end{frame}	
}
	
\begin{frame}[plain,fragile]
	\tiny
\begin{lstlisting}[basicstyle=\tiny]
tekst={}
section={}
length={}
...
...
with open(bestandsnaam) as f:
    for line in f:
        line=line.replace("\r","")
        if line=="\n":
            continue
        matchObj=re.match(r"\s+(\d+) of (\d+) DOCUMENTS",line)
        if matchObj:
            artikelnr= int(matchObj.group(1))
            tekst[artikelnr]=""
            continue
        if line.startswith("SECTION"):
            section[artikelnr]=line.replace("SECTION: ","").rstrip("\n")
        elif line.startswith("LENGTH"):
            length[artikelnr]=line.replace("LENGTH: ","").rstrip("\n")
...
...
...

       else:
           tekst[artikelnr]=tekst[artikelnr]+line
\end{lstlisting}

\end{frame}
	

\begin{frame}[standout]
Let's exercise a bit! 
\end{frame}




\begin{frame}[plain]
    \printbibliography
\end{frame}

\end{document}
